% texte.tex — f169v, vv. 590–634 (conversion initiale des travaux antérieurs)
\vers[n=590]{ΚΛ.}{Clisthène}{ἄφευσεν αὐτὸν κἀπέτιλ᾽ Εὐριπίδης}{Euripide l'a flambé et épilé ;}
\vers[n=591]{}{}{καὶ τἄλλ᾽ ἅπανθ᾽ ὥσπερ γυναῖκ᾽ ἐσκεύασεν.}{et pour tout le reste l'a équipé en femme.}
\vers[n=592]{ΜΝ.}{Mnésiloque}{πείθεσθε τούτῳ ταῦτα; τίς δ᾽ οὕτως ἀνὴρ}{Vous en fiez-vous à lui? Quel homme est assez}
\vers[n=593]{}{}{ἠλίθιος ὅστις τιλλόμενος ἠνείχετο;}{niais pour s'être laissé épiler.}
\vers[n=594]{}{}{οὐκ οἴομαι 'γωγ᾽ ὦ πολυτιμήτω θεώ.}{Je n'en crois rien, pour ma part, ô déesses très honorées.}
\vers[n=595]{ΚΛ.}{Clisthène}{ληρεῖς· ἐγὼ γὰρ οὐκ ἂν ἦλθον ἀγγελῶν,}{Tu veux rire. Je ne serais pas venu l'annoncer}
\vers[n=596]{}{}{εἰ μὴ ᾽πεπύσμην ταῦτα τῶν σάφ᾽ εἰδότων.}{si je ne le tenais de ceux qui le savent bien.}
\vers[n=597]{ΚΟ.}{Coryphée}{τὸ πρᾶγμα τουτὶ δεινὸν εἰσαγγέλλεται.}{C'est une grave affaire qu'on nous annonce là.}
\vers[n=598]{}{}{ἀλλ᾽ ὦ γυναῖκες οὐκ ἐλινύειν ἐχρῆν,}{Or ça, femmes, il s'agit, non de se morfondre,}
\vers[n=599]{}{}{ἀλλὰ σκοπεῖν τὸν ἄνδρα καὶ ζητεῖν ὅπου}{mais de guetter l'homme, de chercher où}
\vers[n=600]{}{}{λέληθεν ἡμᾶς κρυπτὸς ἐγκαθήμενος.}{il se dérobe à nous et se tient caché.}
\vers[n=601]{}{}{καὶ σὺ ξυνέξευρ᾽ αὐτόν, ὡς ἂν τὴν χάριν}{Toi, aide-nous à le trouver, que nous t'ayons}
\vers[n=602]{}{}{ταύτην τε κἀκείνην ἔχῃς ὦ πρόξενε.}{encore cette autre obligation, ô notre patron.}
\vers[n=603]{ΚΛ.}{Clisthène}{φέρ᾽ ἴδω· τίς ἡ πρώτη σύ;}{Voyons, toi d'abord, qui es-tu ?}
\vers[suite]{ΜΝ.}{Mnésiloque}{ποῖ τις τρέψεται;}{(À part :) Où se tourner ?}
\vers[n=604]{ΚΛ.}{Clisthène}{ζητητέαι γάρ ἐστε.}{... Car on va vous passer en revue.}
\vers[suite]{ΜΝ.}{Mnésiloque}{κακοδαίμων ἐγώ.}{(À part :) Malheureux moi !}
\vers[n=605]{ΜΙ.}{Mika}{ἔμ᾽ ἥτις εἴμ᾽ ἤρου; Κλεωνύμου γυνή.}{Tu me demandais qui je suis ? la femme de Cléonyme.}
\vers[n=606]{ΚΛ.}{Clisthène}{γιγνώσκεθ᾽ ὑμεῖς ἥτις ἔσθ᾽ ἡδὶ γυνή;}{Vous savez, vous autres, qui est cette femme ?}
\vers[n=607]{ΚΟ.}{Coryphée}{γιγνώσκομεν δῆτ᾽. ἀλλὰ τὰς ἄλλας ἄθρει.}{Nous le savons. Examine les autres.}
\vers[n=608]{ΚΛ.}{Clisthène}{ἡδὶ δὲ δὴ τίς ἐστιν ἡ τὸ παιδίον}{Et celle-là, qui est-elle ? Celle qui porte}
\vers[n=609]{}{}{ἔχουσα;}{l'enfant ?}
\vers[suite]{ΜΙ.}{Mika}{τίτθη νὴ Δί᾽ ἐμή.}{C'est ma nourrice, par Zeus.}
\vers[suite]{ΜΝ.}{Mnésiloque}{διοίχομαι.}{Je suis perdu. (Il va pour sortir.)}
\vers[n=610]{ΚΛ.}{Clisthène}{αὕτη σὺ ποῖ στρέφει; μέν᾽ αὐτοῦ. τί τὸ κακόν;}{Holà; toi, où vas-tu ? Reste ici. (Mnésiloque s'accroupit.) Qu'est-ce qui t'arrive ?}
\vers[n=611]{ΜΝ.}{Mnésiloque}{ἔασον οὐρῆσαί μ᾽· ἀναίσχυντός τις εἶ.}{Laisse-moi lâcher l'eau. Tu es un impudent.}
\vers[suite]{ΚΛ.}{Clisthène}{σὺ δ᾽ οὖν ποίει τοῦτ᾽. ἀναμενῶ γὰρ ἐνθάδε.}{Eh bien, fais. J'attendrai ici.}
\vers[n=613]{ΚΟ.}{Coryphée}{ἀνάμενε δῆτα καὶ σκόπει γ᾽ αὐτὴν σφόδρα·}{Oui, attends, et surveille-la bien.}
\vers[n=614]{}{}{μόνην γὰρ αὐτὴν ὦνερ οὐ γιγνώσκομεν.}{Elle est la seule que nous ne connaissions pas.}
\vers[n=615]{ΚΛ.}{Clisthène}{πολύν γε χρόνον οὐρεῖς σύ.}{Tu es bien longue à faire de l'eau.}
\vers[suite]{ΜΝ.}{Mnésiloque}{νὴ Δί᾽ ὦ μέλε·}{Par Zeus, l'ami,}
\vers[n=616]{}{}{στραγγουριῶ γάρ· ἐχθὲς ἔφαγον κάρδαμα.}{c'est que j'ai une rétention : hier j'ai mangé du cardame.}
\vers[n=617]{ΚΛ.}{Clisthène}{τί καρδαμίζεις; οὐ βαδιεῖ δεῦρ᾽ ὡς ἐμέ;}{Que parles-tu de cardame ? Veux-tu venir ici, près de moi ? (Il la saisit)}
\vers[n=618]{ΜΝ.}{Mnésiloque}{τί δῆτά μ᾽ ἕλκεις ἀσθενοῦσαν;}{Pourquoi me tirer ainsi ? une valétudinaire !}
\vers[suite]{ΚΛ.}{Clisthène}{εἰπέ μοι,}{Réponds-moi.}
\vers[n=619]{}{}{τίς ἔστ᾽ ἀνήρ σοι;}{Quel est ton mari ?}
\vers[suite]{ΜΝ.}{Mnésiloque}{τὸν ἐμὸν ἄνδρα πυνθάνει;}{Mon mari, dis-tu ?}
\vers[n=620]{}{}{τὸν δεῖνα γιγνώσκεις, τὸν ἐκ Κοθωκιδῶν.}{Connais-tu... chose, du dème de Cothocides.}
\vers[n=621]{ΚΛ.}{Clisthène}{τὸν δεῖνα; ποῖον;}{Chose ? Qui ça ?}
\vers[suite]{ΜΝ.}{Mnésiloque}{ἔσθ᾽ ὁ δεῖν᾽, ὃς καί ποτε}{C'est chose, qui naguère...}
\vers[n=622]{}{}{τὸν δεῖνα τὸν τοῦ δεῖνα—}{avec chose, fils de chose...}
\vers[suite]{ΚΛ.}{Clisthène}{ληρεῖν μοι δοκεῖς.}{Tu m'as l'air de baliverner.}
\vers[n=623]{}{}{ἀνῆλθες ἤδη δεῦρο πρότερον;}{Es-tu déjà venue ici ?}
\vers[suite]{ΜΝ.}{Mnésiloque}{νὴ Δία}{Oui, par Zeus,}
\vers[n=624]{}{}{ὁσέτη γε.}{tous les ans.}
\vers[suite]{ΚΛ.}{Clisthène}{καὶ τίς σοὐστὶ συσκηνήτρια;}{Et quelle est ta compagne de tente ?}
\vers[n=625]{ΜΝ.}{Mnésiloque}{ἡ δεῖν᾽ ἔμοιγ᾽.}{La mienne, c'est chose...}
\vers[suite]{ΚΛ.}{Clisthène}{οἴμοι τάλας. οὐδὲν λέγεις.}{Misère de moi. Tu ne dis rien.}
\vers[n=626]{ΜΙ.}{Mika}{ἄπελθ᾽. ἐγὼ γὰρ βασανιῶ ταύτην καλῶς}{Retire-toi. C'est moi qui vais l'éplucher comme il faut}
\vers[n=627]{}{}{ἐκ τῶν ἱερῶν τῶν πέρυσι· σὺ δ᾽ ἀπόστηθί μοι,}{sur les cérémonies de l'an dernier. — Tiens-toi à distance !}
\vers[n=628]{}{}{ἵνα μὴ 'πακούσῃς ὢν ἀνήρ. σὺ δ᾽ εἰπέ μοι}{tu es un homme, tu ne dois rien entendre. — Toi, dis-moi}
\vers[n=629]{}{}{ὅ τι πρῶτον ἡμῖν τῶν ἱερῶν ἐδείκνυτο.}{lequel des rites sacrés nous fut révélé le premier.}
\vers[n=630]{ΜΝ.}{Mnésiloque}{φέρ᾽ ἴδω, τί μέντοι πρῶτον ἦν; ἐπίνομεν.}{Voyons. Que fût-ce encore que le premier? ... Nous bûmes.}
\vers[n=631]{ΜΙ.}{Mika}{τί δὲ μετὰ τοῦτο δεύτερον;}{Et après ? le second ?}
\vers[suite]{ΜΝ.}{Mnésiloque}{προὐπίνομεν.}{Nous bûmes à nos santés.}
\vers[n=632]{ΜΙ.}{Mika}{ταυτὶ μὲν ἤκουσάς τινος· τί δ᾽ ἦν τρίτον;}{Tu l'as appris de quelqu'une. Et le troisième ?}
\vers[n=633]{ΜΝ.}{Mnésiloque}{σκάφιον Ξένυλλ᾽ ᾔτησεν· οὐ γὰρ ἦν ἀμίς.}{Xenylla demanda une nef, car il n'y avait pas de pot de chambre.}
\vers[n=634]{ΜΙ.}{Mika}{οὐδὲν λέγεις. δεῦρ᾽ ἐλθὲ δεῦρ᾽ ὦ Κλείσθενες·}{Cela ne veut rien dire. Ici, viens ici, ô Clisthène.}
